\section{Bias Temperature Instability}
\label{sec:bg-bti}

Bias temperature instability (BTI) is the dominant transistor wearout
mechanism in CMOS process nodes below 180~nm~\cite{novak2015intel14,
	22nmTrigate, mahapatra2016fundamentals}. BTI manifests as a shift in
transistor threshold voltage ($\Delta V_\text{th}$) under prolonged gate
bias stress. Two complementary mechanisms are distinguished by polarity:
\emph{negative} BTI (NBTI) affects PMOS transistors under negative
$V_\text{GS}$ stress, while \emph{positive} BTI (PBTI) affects NMOS
transistors under positive $V_\text{GS}$
stress~\cite{schroder2003negative, mahapatra2016fundamentals,
	amouri2014aging}. The relative contributions of NBTI and PBTI to
transistor wearout depend strongly on the gate dielectric stack:
high-$\kappa$/metal-gate processes used in nodes below 32~nm exhibit
significant PBTI in NMOS devices, whereas earlier oxynitride-based
processes were dominated by NBTI in
PMOS~\cite{mahapatra2016fundamentals, novak2015intel14, 22nmTrigate}.
The defects underlying both effects are believed to involve a
combination of hydrogen-related traps and charge carriers in the gate
dielectric~\cite{bti_modelling_2013, mahapatra2016fundamentals}.

From the side-channel perspective, the feature that distinguishes BTI
from other wearout mechanisms such as hot-carrier injection (HCI) and
time-dependent dielectric breakdown (TDDB) is its partial
recoverability. When the stressing bias is removed or reversed, $\Delta
	V_\text{th}$ recovers toward its pre-stress value on time scales ranging
from milliseconds to hours, depending on temperature and stress
conditions~\cite{stott2010degradation, ramey2014bti, grasser2016PRNBTI,
	bti_modelling_2013}. This stress-recovery asymmetry produces an
observable signature of a transistor's recent bias history: a transistor
that has been recently stressed exhibits a temporarily-elevated
$V_\text{th}$ that is detectably distinct from one in its rested state.
BTI also includes a permanent (non-recoverable) component that
accumulates with cycling and saturates around 30~mV after $\sim 10^4$
seconds of cumulative stress~\cite{grasser2011permanentifc,
	grasser2016PRNBTI}. From a process-engineering perspective, the
recoverable component dominates short-term reliability concerns, while
the permanent component contributes to slow drift in transistor
parameters over the device lifetime.

Both reaction-diffusion (R-D) model and the more recent stochastic
capture-emission framework~\cite{bti_modelling_2013,
	mahapatra2016fundamentals} predict that $\Delta V_\text{th}$ follows a
power law in cumulative stress time:
\begin{equation}
	\label{eqn:bti-powerlaw}
	\Delta V_\text{th} = b \cdot \alpha^n \cdot t^n,
\end{equation}
where $b$ is a technology- and temperature-dependent constant,
$\alpha$ is the duty cycle of bias stress, $n$ is a
process-dependent BTI exponent typically in the range
$0.1$--$0.25$, and $t$ is cumulative stress time~\cite{wang2007bti,
	nunes2013ffbti, bti_modelling_2013}. $\Delta V_\text{th}$ propagates to a measurable shift in transistor
switching delay. For an inverter, the dependence of propagation
delay on accumulated BTI can be approximated as
\begin{equation}
	\label{eqn:bti-tp}
	t_p = a_0 + a_1 \cdot \alpha^n \cdot t^n,
\end{equation}
where $a_0$ is the intrinsic gate delay and $a_1$ is a fitting
parameter that depends on transistor sizing and the process node's
sensitivity to $\Delta V_\text{th}$~\cite{wang2007bti,
	nunes2013ffbti}. Process engineers typically target
$\Delta V_\text{th} \leq 0.05 \cdot V_\text{DD}$ at five years of
stress~\cite{taek2013bti, rzepa2018comphy}; for nodes below 32~nm
this corresponds to propagation-delay shifts on the order of single
picoseconds, comparable to the resolution of fine-grained
time-to-digital converters implemented in FPGA soft logic
(Section~\ref{sec:bg-tdc}).

FPGAs are particularly well-suited to BTI measurement because their
reconfigurable fabric permits arbitrary placement of fine-grained timing
instrumentation that directly connects to the routing nets being
probed~\cite{stott2010degradation, amouri2014aging,
	drewes2024pentimento}. When a single bit of information is held in a
routing net for an extended period, the aggregate BTI stress across all
transistors carrying that bit accumulates, producing a propagation-delay
shift detectable through on-chip timing measurement. Drewes et
al.~\cite{drewes2024pentimento} demonstrated this principle as a
temporal-multi-tenancy side-channel: an adversary co-located on a target
FPGA after a victim's workload has completed can recover bits of
information stored in routing nets by the prior tenant, even after the
FPGA has been reconfigured. The attack relies on (i) BTI's recoverable
component producing data-dependent delay shifts and (ii) the FPGA's TDC
infrastructure resolving those shifts at picosecond granularity. Whether
analogous defect dynamics in other device families, and in particular
the filament-formation dynamics of ReRAM (Sections~\ref{sec:bg-oxram}
and~\ref{sec:bg-defects}), support similar attacks is an open question
that motivates the remainder of this dissertation.